\documentclass[11pt]{article}

\usepackage[utf8]{inputenc}
\usepackage[T1]{fontenc}
\usepackage{mathptmx}
\usepackage[margin=1.0in]{geometry}
\usepackage{amsmath,amssymb}
\usepackage{xcolor}
\usepackage{parskip}
\usepackage{enumitem}
\usepackage[colorlinks=true,linkcolor=blue!60!black,urlcolor=blue!60!black]{hyperref}

% ==========================================================================
% AUTO-GENERATED FILE — DO NOT EDIT MANUALLY
% Generated by: manuscript/scripts/generate_latex_stats.py
% Source data:  manuscript/analysis_results.json
% ==========================================================================

% --- Participant & Sample Stats ---
\newcommand{\statNRecruited}{35}
\newcommand{\statNExcludedEOG}{7}
\newcommand{\statNIncluded}{28}
\newcommand{\statAgeRange}{17--34}
\newcommand{\statEOGRejectPct}{30\%}
\newcommand{\statHEOGStepUV}{25\,\mu\text{V}}
\newcommand{\statCompRate}{\$10/\text{hr}}

% --- Task & Staircase Parameters ---
\newcommand{\statViewDistCm}{57\,\text{cm}}
\newcommand{\statRefreshRateHz}{120\,\text{Hz}}
\newcommand{\statTargetDurMs}{8.33\,\text{ms}}
\newcommand{\statMaskDurMs}{8.33\,\text{ms}}
\newcommand{\statTargetMaskSOAMs}{41.7\,\text{ms}}
\newcommand{\statPostMaskBlankMs}{516.6\,\text{ms}}
\newcommand{\statFixationDurMs}{700\,\text{ms}}
\newcommand{\statCueDursMs}{1242, 1284, 1325, or 1367\,\text{ms}}
\newcommand{\statNOrientations}{24}
\newcommand{\statNBlocks}{8}
\newcommand{\statTrialsPerBlock}{48}
\newcommand{\statPracticeTrials}{20}
\newcommand{\statTotalTrials}{404}
\newcommand{\statStaircaseRate}{0.65}
\newcommand{\statStaircaseBlocks}{3}
\newcommand{\statStaircaseGrayMean}{54.0}
\newcommand{\statStaircaseGraySD}{22.0}

% --- Behavioral ANOVA (Fmax) Statistics ---
\newcommand{\statBehavDF}{1,27}
\newcommand{\statRightSideFloorDeg}{28^\circ}
\newcommand{\statAbsErrSideFmax}{11.32}
\newcommand{\statAbsErrSideCrit}{\pm 3.95}
\newcommand{\statAbsErrSideP}{0.001}
\newcommand{\statAbsErrCueFmax}{12.30}
\newcommand{\statAbsErrCueCrit}{\pm 3.96}
\newcommand{\statAbsErrCueP}{< 0.001}
\newcommand{\statAbsErrIntFmax}{9.31}
\newcommand{\statAbsErrIntCrit}{\pm 3.85}
\newcommand{\statAbsErrIntP}{0.002}
\newcommand{\statRelErrSideFmax}{2.04}
\newcommand{\statRelErrSideCrit}{\pm 4.09}
\newcommand{\statRelErrSideP}{0.169}
\newcommand{\statRelErrCueFmax}{< 1}
\newcommand{\statRelErrIntFmax}{1.94}
\newcommand{\statRelErrIntCrit}{\pm 4.14}
\newcommand{\statRelErrIntP}{0.176}

% --- FOOOF Spectral Parameterization Statistics ---
\newcommand{\statFOOOFFreqRange}{3.9 to 48.8\,\text{Hz}}
\newcommand{\statFOOOFFreqRes}{1.9\,\text{Hz}}
\newcommand{\statFOOOFPeakWidth}{[1, 5]\,\text{Hz}}
\newcommand{\statFOOOFMaxPeaks}{7}
\newcommand{\statFOOOFMinHeight}{0.22}
\newcommand{\statFOOOFPeakThresh}{0.20}
\newcommand{\statFOOOFBonfAlpha}{0.0125}
\newcommand{\statFOOOFOffsetCrit}{\pm 15.40}
\newcommand{\statFOOOFExpCrit}{\pm 15.35}
\newcommand{\statFOOOFRsqCrit}{\pm 14.94}
\newcommand{\statFOOOFErrCrit}{\pm 16.76}
\newcommand{\statIAFGrandMeanHz}{10.6\,\text{Hz}}
\newcommand{\statIAFGrandSDHz}{0.5\,\text{Hz}}
\newcommand{\statIAFRangeHz}{9.6--11.4\,\text{Hz}}

% --- MVPA & Permutation Parameters ---
\newcommand{\statMVPAPerms}{10,000}
\newcommand{\statMVPAAlpha}{0.05}
\newcommand{\statMVPABonfAlpha}{0.005}
\newcommand{\statCVFolds}{10}
\newcommand{\statCVRepeatsEmp}{10}
\newcommand{\statCVTotalEmp}{100}
\newcommand{\statCVRepeatsNull}{20}
\newcommand{\statCVTotalNull}{200}
\newcommand{\statSVMCostC}{1}
\newcommand{\statSVRCostC}{0.1}
\newcommand{\statAlphaWinMs}{20\,\text{ms}}
\newcommand{\statAlphaStepMs}{20\,\text{ms}}
\newcommand{\statAlphaNPoints}{50}
\newcommand{\statCFCSlidingWinMs}{250\,\text{ms}}
\newcommand{\statCFCWinMs}{15\,\text{ms}}
\newcommand{\statCFCStepMs}{15\,\text{ms}}
\newcommand{\statCFCNPoints}{50}
\newcommand{\statFCSlidingWinMs}{375\,\text{ms}}
\newcommand{\statFCWinMs}{50\,\text{ms}}
\newcommand{\statFCStepMs}{50\,\text{ms}}
\newcommand{\statFCNPoints}{12}
\newcommand{\statFCZThreshWin}{0.45}
\newcommand{\statFCZThreshEpoch}{0.30}
\newcommand{\statSVMAlphaPctSigRange}{90\%--96\%}
\newcommand{\statSVMAlphaPctSigMin}{90\%}
\newcommand{\statSVMAlphaLeftWin}{-1000 to -820 ms and -760 to -20 ms}
\newcommand{\statSVMAlphaRightWin}{-1000 to -800 ms, -760 to -460 ms, -420 to -200 ms, -160 to -80 ms, and -20 ms}
\newcommand{\statSVMMVLPctSigRange}{8\%--14\%}
\newcommand{\statSVMMVLLeftSigMs}{-875\,\text{ms}}
\newcommand{\statSVMMVLRightSigWinMs}{-785 to -770\,\text{ms}}
\newcommand{\statSVMMVLRightSigPtMs}{-380\,\text{ms}}
\newcommand{\statSVMwPLIPctSigRange}{0\%--8.3\%}
\newcommand{\statSVMwPLIRightSigMs}{-462.5\,\text{ms}}
\newcommand{\statSVRAlphaRightSigMs}{-640\,\text{ms}}
\newcommand{\statSVRAACLeftSigMs}{-200\,\text{ms}}
\newcommand{\statSVRCorrLeftSigMs}{-562.5\,\text{ms}}

% --- Dynamic Figure File Paths ---
\newcommand{\figPathOne}{figures/figure1_task.png}
\newcommand{\figPathTwo}{figures/figure2_behavior.png}
\newcommand{\figPathThree}{figures/figure3_svm_cfc.png}
\newcommand{\figPathFour}{figures/figure4_svm_fc.png}
\newcommand{\figPathFive}{figures/figure5_svr_cfc.png}
\newcommand{\figPathSix}{figures/figure6_svr_fc.png}
\newcommand{\figPathSeven}{figures/figure7_mvpa_summary.png}


\definecolor{revblue}{RGB}{0, 75, 145}
\definecolor{commentgray}{RGB}{70, 70, 70}
\newcommand{\revcomment}[1]{\vspace{6pt}\noindent\textbf{\textcolor{commentgray}{#1}}\par\vspace{2pt}}
\newcommand{\authorresp}[1]{\noindent\textcolor{revblue}{\textbf{Response:}} #1\par\vspace{6pt}}

\begin{document}

\begin{center}
{\Large \textbf{Response to Decision Letter and Reviewer Comments}}\\[6pt]
{\large \textbf{Manuscript ID: CerCor-2022-00354}}\\[4pt]
\textit{Connecting Covert Attention and Visual Perception to the Spatiotemporal Dynamics of Alpha Band Activity, Cross-Frequency Coupling (CFC), and Functional Connectivity using Multivariate Pattern Analysis (MVPA)}\\
Sarah S.\ Sheldon, Alona Fyshe, and Kyle E.\ Mathewson
\end{center}

\vspace{8pt}
\hrule
\vspace{8pt}

We thank the Editor-in-Chief and both reviewers for their thorough reading of our manuscript and their constructive suggestions. Below, we provide a point-by-point response to each comment and indicate the corresponding revisions in the manuscript (where new and revised sections are highlighted in blue). We have also complied with all editorial formatting requirements: the revised manuscript is compiled as a single double-spaced PDF with embedded tables and figures at their relevant locations, and all references are formatted strictly in \textit{Cerebral Cortex} style.

% ==============================================================================
% REVIEWER 1
% ==============================================================================
\section*{Reviewer 1}

\revcomment{1. The authors acknowledge in the discussion that using EEG to measure gamma-band activity comes with limitations. It is also sub-optimal for functional connectivity metrics. This being the case, the choice of EEG to test the stated hypotheses is not optimal. I wonder whether the authors might be able to explicitly state, perhaps in the discussion, exactly what the current results can confidently add to the literature given these limitations? The fact that pre-stimulus alpha amplitude dissociates between cued and non-cued trials is certainly not a novel finding, and I have concerns about how much one should read into the CFC and connectivity results given the inherent limitations of using EEG to capture these measures (see further points below for details). The analysis approach is sophisticated, but I found myself struggling to grasp exactly what the study adds that wasn't known already.}

\authorresp{We appreciate this important point and have added a dedicated subsection in the Discussion (\textit{``What the Current MVPA Comparison Adds to the Literature and Signal-to-Noise Considerations''}) as well as a new multi-panel summary comparison figure (\textbf{Figure 7}) to make our core contribution explicit. Specifically, while univariate reductions in posterior alpha amplitude following spatial cues are well known, theoretical models such as the oscillation-based probability of response (OPR) model (Zazio et al.\ 2020) and network synchronization accounts (Sadaghiani and Kleinschmidt 2016) propose that macroscopic alpha amplitude changes are secondary to local alpha--gamma phase--amplitude coupling (PAC) or inter-areal alpha phase synchronization. To our knowledge, no prior study had benchmarked single-trial spatiotemporal alpha amplitude head-to-head against five local CFC estimators and four volume-conduction-mitigated functional connectivity estimators within the exact same participants ($N = \statNIncluded{}$), task epochs, and cross-validated decoding pipeline for both attentional classification (SVM) and continuous perceptual error prediction (SVR). Our results establish that at the scalp EEG level, spatiotemporal alpha amplitude carries substantially more robust and temporally sustained information about covert spatial attention ($\statSVMAlphaPctSigRange{}$ of pre-stimulus windows) than local low-gamma CFC or pairwise sensor-level alpha connectivity, while simultaneously demonstrating that this strong attentional signal does not translate into trial-by-trial prediction of continuous orientation reproduction accuracy---directly supporting the baseline sensory excitability model (BSEM; Samaha et al.\ 2020).}

\revcomment{2. One issue that I wonder about is whether differences in the classification performance of the different EEG measures may be partly attributable to differences in their signal-to-noise ratios? For instance, alpha amplitude is likely a less noisy estimate which is easier to measure with EEG than cross frequency coupling metrics. Might differences in signal-to-noise represent a confounding factor for interpreting the results?}

\authorresp{We completely agree and now explicitly discuss single-trial signal-to-noise ratio (SNR) as a key methodological consideration in the Discussion (\textit{``What the Current MVPA Comparison Adds to the Literature and Signal-to-Noise Considerations''}). As we now explain, instantaneous alpha amplitude is a first-order narrowband envelope estimate of the dominant human scalp rhythm and therefore enjoys a high single-trial SNR. By contrast, CFC and functional connectivity are higher-order relational statistics that require estimating fast, low-amplitude gamma envelopes ($30\text{--}45\text{ Hz}$) or pairwise inter-electrode phase relationships over short sliding windows ($250\text{--}375\text{ ms}$), which inherently exhibit higher estimation variance and lower single-trial SNR at the scalp (Aru et al.\ 2015; Bastos and Schoffelen 2016; H\"ulsemann et al.\ 2019). We explicitly note that the superior decoding performance of alpha amplitude reflects both its role in modulating baseline sensory excitability and its higher single-trial estimator SNR in non-invasive scalp EEG.}

\revcomment{3. The task employed was quite different to tasks usually adopted by studies investigating the influence of prestimulus alpha on perception (i.e. usually simple hit/miss detection tasks of 2-AFC discrimination tasks). One potentially important difference is that the task employed here required participants to use the computer mouse to rotate a needle to reproduce the direction of the presented target. It seems that errors on this task will be influenced by more factors than just visual perceptual factors (i.e. motor errors will potentially also have a strong influence). Do the authors think that this might have contributed to the lack of EEG-behavioural effects that they found?}

\authorresp{Yes, this is an insightful observation. We have added a dedicated discussion of this point in the new Discussion subsection \textit{``Task Demands, Continuous Motor Reproduction, and Visual Field Asymmetry''}. We contrast our continuous $360^\circ$ mouse-rotation reproduction task (which includes a $\statPostMaskBlankMs{}$ post-mask delay followed by self-paced continuous mouse rotation) against standard binary yes/no detection and 2-AFC discrimination tasks. As we now discuss, single-trial continuous orientation response errors incorporate variance not only from early visual encoding, but also from visual working memory maintenance across the post-mask delay, visuomotor coordinate transformation, and fine motor execution with the computer mouse. Because pre-stimulus oscillations primarily modulate early sensory excitability around target onset, downstream working memory and motor execution variance likely attenuated the trial-by-trial correlation ($r_z$) between pre-stimulus EEG patterns and final orientation reproduction error.}

\revcomment{4. The following is from the Methods section regarding multiple comparison correction: ``For group level statistical analysis, t-tests using a threshold of p < 0.05 were used to compare the empirical results with permuted-labels results (Bode et al. 2019). Subsequently, right-tailed tmax permutation testing (number of iterations = 10000, alpha level = 0.05) was conducted for multiple comparisons correction (Blair and Karniski 1993).'' My interpretation is that the multiple comparison correction was applied within each separate analysis (across timepoints?). However, given the number of different EEG metrics being tested (and for left and right targets separately etc.), are the authors confident that the type-1 error rate is being controlled to an acceptable degree?}

\authorresp{Your interpretation is correct: right-tailed $t_{\max}$ permutation testing ($\statMVPAPerms{}$ iterations, $\alpha = \statMVPAAlpha{}$) controlled the familywise error rate across all time points within each separate metric and target-side analysis. We have clarified this in both the \textit{Multivariate Pattern Analysis (MVPA)} Methods section and the Discussion. Specifically, we explain that we retained within-analysis $t_{\max}$ correction so as not to overly penalize the statistical power of the CFC and functional connectivity analyses relative to alpha amplitude. Crucially, however, we now explicitly state that isolated transient significance spikes (such as the sparse SVM windows in $\text{MVL}_{\text{\"Ozkurt}}$ [$\statSVMMVLLeftWin{}$ for left targets; $\statSVMMVLRightWin{}$ for right targets], MI [$\statSVMMILeftWin{}$ for left targets; $\statSVMMIRightWin{}$ for right targets], and $\text{PAC}_{\text{GLM}}$ [$\statSVMGLMRightWin{}$ for right targets], or the single-window SVR effects in AAC [$\statSVRAACRightWin{}$], $\text{MVL}_{\text{\"Ozkurt}}$ [$\statSVRMVLRightWin{}$], and $\text{PAC}_{\text{GLM}}$ [$\statSVRGLMLeftWin{}$]) would \textit{not} survive an additional Bonferroni correction across the 10 tested metrics ($\alpha_{\text{Bonf}} = \statMVPABonfAlpha{}$) and should therefore be interpreted with caution, whereas the sustained pre-stimulus decoding of alpha amplitude ($\statSVMAlphaPctSigRange{}$ of all pre-stimulus windows, $p < 0.001$) remains robust under any experiment-wise correction threshold.}

\revcomment{5. There seems to be a visual field asymmetry in behavioural performance (Figure 2), with larger errors on the left than the right side. Could the authors comment on why this might be?}

\authorresp{We have added a detailed explanation of this visual field asymmetry in the Discussion subsection \textit{``Task Demands, Continuous Motor Reproduction, and Visual Field Asymmetry''}. Specifically, absolute response errors were significantly larger for left-hemifield targets than right-hemifield targets (Figure~2A; main effect of Target Side: $F_{\max}(\statBehavDF{}) = \statAbsErrSideFmax{}, p \statAbsErrSideP{}$), and informative spatial cues significantly reduced error magnitude on the left side while right-side error was already low ($\sim \statRightSideFloorDeg{}$) on non-informative trials. We discuss two complementary factors: (1) all $\statNIncluded{}$ participants were right-handed and used their right hand to control the mouse during continuous visuomotor rotation, which favors right-visual-field (left-hemisphere) sensorimotor integration; and (2) the preliminary target-detection staircase adjusted target contrast pooled across both visual hemifields rather than separately per hemifield, placing right-hemifield targets closer to floor error even without an informative spatial cue and leaving greater dynamic range for covert attention to reduce errors on left-hemifield targets.}

% ==============================================================================
% REVIEWER 2
% ==============================================================================
\section*{Reviewer 2}

\revcomment{General: The general approach, implementation and results of this manuscript seem solid. However, the methods are not always described in a way that makes it clear precisely what was done. In a few spots, the writing is too colloquial. Further, some passages especially in the discussion are confusing and do not seem to reflect a thorough understanding of electrophysiology. The results nicely fit with the hypotheses. If possible, it might be nice to have a ``summary'' figure where the effectiveness of the different metrics in MVPA was compared. Figure 3-6 are hardly readable.}

\authorresp{We thank Reviewer 2 for these constructive suggestions. In response, we have: (1) created and embedded a new 3-panel summary comparison figure (\textbf{Figure 7}) comparing peak SVM classification accuracy, percentage of significant pre-stimulus decoding windows, and peak SVR prediction accuracy across all 10 EEG metrics; (2) revised all colloquial phrasing and clarified every methodological detail noted below; (3) rewritten the electrophysiology discussion comparing ECoG, MEG, and scalp EEG high-gamma recordings; and (4) configured our automated figure generation pipeline to enforce high-resolution readability, abbreviated axis labels (\texttt{MI}, \texttt{Accuracy}, $r_z$), and $200\text{ ms}$ tick spacing.}

\subsection*{Abstract}
\begin{itemize}[leftmargin=1.5em]
    \item \textbf{``The listing of 3 different possible functions for alpha is too elaborate for the abstract''} --- \textcolor{revblue}{We condensed the opening of the Abstract into a single concise sentence introducing the three candidate mechanisms (local alpha amplitude, local alpha--gamma CFC, and inter-electrode alpha functional connectivity).}
    \item \textbf{``Abbreviation CFC needs to be explained in the abstract''} --- \textcolor{revblue}{``cross-frequency coupling (CFC)'' is now spelled out at first mention in the first sentence of the Abstract.}
    \item \textbf{``\,'pre-target' should perhaps be renamed 'pre-stimulus' in the abstract''} --- \textcolor{revblue}{Changed to ``pre-stimulus'' throughout the Abstract and harmonized across the manuscript.}
    \item \textbf{``Metrics cannot have activity, they can show time courses or patterns''} --- \textcolor{revblue}{Revised throughout the Abstract and main text to refer to ``time courses'' or ``spatiotemporal patterns'' of each metric rather than ``activity of each metric.''}
    \item \textbf{``The terms 'SVM' and 'SVR' should be explained. Or maybe do not refer to these specific methods at all in the abstract.''} --- \textcolor{revblue}{Both ``linear support vector machine (SVM) classification'' and ``linear support vector regression (SVR)'' are now fully spelled out at first mention in the Abstract.}
    \item \textbf{``\,'alpha amplitude alone is most sensitive' -- this sounds as if alpha amplitude alone was better than when combined with CFC and/or synchrony, which was not tested for.''} --- \textcolor{revblue}{Rephrased to: \textit{``Among the tested metrics, pre-stimulus alpha amplitude was most sensitive to the deployment of covert spatial attention.''}}
\end{itemize}

\subsection*{Introduction}
\begin{itemize}[leftmargin=1.5em]
    \item \textbf{``Background and own approach/hypothesis should be more separate in the introduction... Would be good to have a clearer structure here.''} --- \textcolor{revblue}{We restructured the Introduction so that the first four paragraphs present the theoretical background (BSEM, OPR model, and inter-areal alpha functional connectivity, plus how they relate to one another), and the final paragraph cleanly introduces our experimental approach, MVPA methodology, and specific testable predictions (summarized in Table 1).}
    \item \textbf{``P6, l.117-9: this sentence should not end in a question mark.''} --- \textcolor{revblue}{Corrected to a declarative sentence with a period.}
    \item \textbf{``P6, l.122: 'the idea is' is too colloquial for a scientific manuscript''} --- \textcolor{revblue}{Removed and replaced with formal scientific phrasing (\textit{``In the current study, we directly compared the information content of...''}).}
\end{itemize}

\subsection*{Methods}
\begin{itemize}[leftmargin=1.5em]
    \item \textbf{``P8, l.164: what is meant by 'identical to the fixation'?''} --- \textcolor{revblue}{Clarified to read: \textit{``Following mask offset, a \statPostMaskBlankMs{} post-mask delay interval occurred during which the display returned to the initial fixation configuration (the central white fixation dot and the two empty dark gray placeholder circles).''}}
    \item \textbf{``P9, l.176: 'rest at their own pace' means 'take a break for as long as they wanted'?''} --- \textcolor{revblue}{Clarified to read: \textit{``Between blocks (every \statTrialsPerBlock{} trials), participants were given a self-paced break and could rest for as long as they wished before initiating the next block.''}}
    \item \textbf{``P11, l.230: Please explain what is the difference between (relative?) 'response error' and 'absolute response error'?''} --- \textcolor{revblue}{We added an explicit mathematical and conceptual definition of both \textit{signed response error} (preserving clockwise/counterclockwise direction in $[-180^\circ, +180^\circ]$ to test for directional rotational bias) and \textit{absolute response error} ($|\text{Target} - \text{Response}| \in [0^\circ, 180^\circ]$, measuring non-negative error magnitude/precision).}
    \item \textbf{``P11, l.233: please formulate better''} --- \textcolor{revblue}{Rewritten formally: \textit{``Compared with parametric ANOVAs, permutation-based $F_{\max}$ testing requires fewer distributional assumptions and provides strong control of familywise Type I error rates with minimal loss of statistical power (Fields and Kuperberg 2019).''}}
    \item \textbf{``P13, l.263-5: how do you come to the conclusion that they 'did not use covert attention as instructed'?''} --- \textcolor{revblue}{Clarified to explain precisely that losing $>\statEOGRejectPct{}$ of trials to the $\statHEOGStepUV{}$ split-half sliding-window HEOG step criterion indicates that those participants frequently made overt horizontal eye movements (saccades) during the cue--target interval rather than maintaining central fixation, leaving too few artifact-free covert-attention trials for \statCVFolds{}-fold MVPA cross-validation.}
    \item \textbf{``P13, l.268ff: technically, the FOOOF only separates 1/f component from peak components, which does not necessarily equate to aperiodic and periodic components, respectively. However, for the purposes here (in alpha band which has a strong peak), this method is sufficient.''} --- \textcolor{revblue}{We updated the subsection title and text (\textit{``Spectral Parameterization of 1/f and Peak Alpha Components''}) to explicitly note that FOOOF mathematically separates the broadband $1/f$-like background decay from narrowband Gaussian spectral peaks, which effectively isolates the prominent alpha peak.}
    \item \textbf{``P14, l.296-7: Are these individual alpha frequencies reported in the manuscript? Were they used in analyses?''} --- \textcolor{revblue}{We now explicitly report the grand-average median individual alpha peak frequency across participants, electrodes, and conditions ($M = \statIAFGrandMeanHz{}, \text{SD} = \statIAFGrandSDHz{}$, condition medians spanning $\statIAFRangeHz{}$) and clarify how each participant's electrode- and condition-specific median peak alpha frequency centered the $\pm 1\text{ Hz}$ bandpass filters for both the instantaneous alpha amplitude and local CFC analyses.}
    \item \textbf{``P15-16: How did you decide that other types of CFC (notably phase-phase coupling aka cross-frequency phase synchrony) are not relevant for the study? Also, please write that you only estimated local CFC (as opposed to inter-areal CFC). Please explain the filtering for the high frequency better. Did you extract a single signal for a whole frequency range? And was this 15-60 or 30-45 Hz? (the former would include beta band)''} --- \textcolor{revblue}{We thoroughly revised the \textit{Local Cross-Frequency Coupling (CFC)} subsection to address all three points: (1) we explicitly state that we estimated \textit{local} (within-electrode) CFC rather than inter-areal CFC; (2) we explain that AAC and PAC were selected over phase--phase ($n:m$) coupling because AAC and PAC are the two specific biophysical mechanisms posited by the OPR model (Zazio et al.\ 2020; Chaumon and Busch 2014); and (3) we clarify that a single broadband high-frequency analytic signal was extracted per trial centered on the $30\text{--}45\text{ Hz}$ low-gamma band using a wide sideband-preserving filter ($\pm 0.4 \times$ center frequency around $30\text{--}45\text{ Hz}$, resulting in a $15\text{--}60\text{ Hz}$ passband bounded at $50\text{ Hz}$ by the preprocessing low-pass filter, following Aru et al.\ 2015 and Seymour et al.\ 2017) because capturing alpha-rate ($8\text{--}14\text{ Hz}$) modulations of the gamma amplitude envelope mathematically requires a filter bandwidth at least twice the modulating alpha frequency, while noting that the lower passband transition extends into the upper beta range.}
\end{itemize}

\subsection*{Results \& Figures}
\begin{itemize}[leftmargin=1.5em]
    \item \textbf{``Please avoid 'As can be seen in...' formulations''} --- \textcolor{revblue}{All instances of ``As can be seen in...'' and ``As shown in...'' have been removed across the Results section.}
    \item \textbf{``Figure 2 should be divided into A and B panel''} --- \textcolor{revblue}{Figure 2 is now divided into explicit bold \textbf{A}, \textbf{B}, and \textbf{C} panels (with significance brackets on Panel A), and the caption and text reference Figure 2A, Figure 2B, and Figure 2C directly.}
    \item \textbf{``P25, l.550-7: Why is the response error mentioned first in text, but shown second in the figure?''} --- \textcolor{revblue}{We reordered the Behavioral Results text so that \textit{absolute response error} (Figure 2A) is presented first, followed by \textit{signed response error} (Figure 2B), matching the visual left-to-right order of Figure 2.}
    \item \textbf{``If possible, it might be nice to have a 'summary' figure where the effectiveness of the different metrics in MVPA was compared.''} --- \textcolor{revblue}{We added \textbf{Figure 7}, a 3-panel summary figure comparing peak SVM classification accuracy (Panel A), percentage of significant pre-stimulus decoding windows (Panel B), and peak SVR prediction correlation $r_z$ (Panel C) across all 10 EEG metrics.}
    \item \textbf{``Figures 3-6: Currently, these figures have considerably too low resolution, too small sub-panels, and too small font size to be readable. Some abbreviations or shortenings may be used in labels if they are explained in legend (e.g. using just MI instead of Modulation Index; using just Accuracy instead of Classification Accuracy). Time points could be given only every 200 ms, but in larger font.''} --- \textcolor{revblue}{We reconstructed Figures 3--7 at 300 DPI directly from the raw data in our end-to-end reproducibility pipeline (\texttt{run\_mvpa\_and\_figures.py} and \texttt{pipeline\_api.py} at \url{https://github.com/APPLabUofA/OrientExoML}), enforcing $200\text{ ms}$ time ticks ($-1000, -800, -600, -400, -200, 0\text{ ms}$) in enlarged typography and concise axis/row labels (\texttt{AAC}, \texttt{Alpha Amplitude}, \texttt{MI}, $\text{MVL}_{\text{\"Ozkurt}}$, $\text{PAC}_{\text{GLM}}$, $\text{PAC}_{\text{PLV}}$, \texttt{Correlation}, \texttt{wPLI}, \texttt{PLV}, \texttt{iCOH}) explained in the figure legends.}
\end{itemize}

\subsection*{Discussion}
\begin{itemize}[leftmargin=1.5em]
    \item \textbf{``P29, l.628: Shouldn't this be 'in/during the 1000 ms before target onset'? The way it is written now, might be interpreted as 'exactly 1000 ms before'\,''} --- \textcolor{revblue}{Revised to \textit{``throughout the $1000\text{ ms}$ interval prior to target onset''}.}
    \item \textbf{``P29, l.632: I recommend writing 'which predicts that changes'\,''} --- \textcolor{revblue}{Updated as recommended.}
    \item \textbf{``P30, l.663-74: Which points are you trying to make here? You are comparing to Voytek's finding of theta-high-gamma PAC involving frequency ranges that were not investigated in this study. Also, what is meant by surface EEG here? Why should it not be possible to record high-gamma from outside the head (which is what Bonnefond \& Jensen 2015 did with MEG)?''} --- \textcolor{revblue}{We completely rewrote this paragraph (\textit{``Implications for the OPR Model and Electrophysiological Scope of Scalp EEG''}). Specifically: (1) we clarify that Voytek (2010) observed a task-dependent shift from theta--high-gamma to alpha--high-gamma ($80\text{--}150\text{ Hz}$) PAC over posterior visual cortex using invasive subdural ECoG; (2) we replace the ambiguous term ``surface EEG'' with ``non-invasive scalp EEG''; and (3) we explicitly acknowledge that high-gamma ($80\text{--}120\text{ Hz}$) PAC \textit{can} be recorded non-invasively from outside the head using MEG, as demonstrated by Bonnefond and Jensen (2015), whereas scalp EEG is more heavily affected by cranial spatial smearing and scalp muscle artifacts above $50\text{ Hz}$ (and our analysis was restricted to low-gamma [$30\text{--}45\text{ Hz}$] below our $50\text{ Hz}$ low-pass filter).}
\end{itemize}

\end{document}
